\section{Limitations and Failure Modes}

\subsection{Calibration Dependency}
The efficacy of \texttt{proxsparse\_pruner} depends on representative calibration sequences. Distributional shift between calibration tokens and deployment data can introduce representation degradation.

\subsection{Hardware and Architectural Constraints}
\begin{itemize}
\item \textbf{Rank Granularity:} Factorized ranks must divide cleanly by attention head count ($r \pmod H == 0$) to avoid irregular tensor padding.
\item \textbf{Kernel Launch Overhead:} On short sequences ($N < 16$), executing dual factorized matrix products incurs launch overhead compared to a single fused GEMM.
\item \textbf{Hardware Acceleration Boundaries:} Hardware speedup for 2:4 structured sparsity requires tensor cores supported on Ampere or newer architectures.
\end{itemize}

\subsection{Numerical Precision Limits}
Aggressive low-rank truncation can cause floating-point underflow under 16-bit precision, requiring explicit safety clamping ($10^{-4}$) to prevent numerical divergence.
